% !PRINTED 441
\weberpage{441}
%<<CONT
vanish. Such a straight line we shall call a \gesperrt{principal axis of the substitution} $A$. But only certain special substitutions possess principal axes. A substitution with a principal axis has, besides this, also \gesperrt{a pole}, which in special cases may also fall into the principal axis.

In order to survey these relations, let us place the principal axis in the line $x_1=0$; then, if $x_1=0$, we must have $y_1=0$, $y_2:y_3=x_2:x_3$. From this it follows that
\[
\begin{aligned}
y_1 &= \alpha\, x_1\\
y_2 &= \alpha'\, x_1 + \beta\, x_2\\
y_3 &= \alpha''\, x_1 + \beta\, x_3.
\end{aligned}
\]
The pole existing outside the principal axis we obtain by setting $y_1=\alpha\,x_1$, $y_2=\alpha\,x_2$, $y_3=\alpha\,x_3$; then we obtain
\[
(\beta-\alpha)\,x_2+\alpha'\,x_1=0,\quad (\beta-\alpha)\,x_3+\alpha''\,x_1=0,
\]
and this point will fall into the line $x_1=0$ when, and only when, $\beta=\alpha$. This case cannot occur with finite groups (except in the case of multiplication).

If a pole distinct from the principal axis exists, we can place it at the vertex $x_2=0$, $x_3=0$, and the substitution takes the form
\[
y_1=\alpha\,x_1,\quad y_2=\beta\,x_2,\quad y_3=\beta\,x_3,
\]
where $\alpha$ is different from $\beta$.

A straight line $u_1x_1+u_2x_2+u_3x_3=0$ is its own image when, and only when, either $u_2=u_3=0$ or $u_1=0$.

\gesperrt{If, then, a substitution has a pole and a principal axis, there remain unchanged, besides the latter, all the straight lines, and only those, which pass through the pole.}

It should further be remarked that, if three separate poles are present, they cannot fall on a straight line without the line becoming a principal axis. For if a straight line is its own image, then, unless the line corresponds to itself point for point, only two points can lie on it which are their own images.

Summing all this up, we find the following kinds of ternary linear substitutions, where coinciding poles are counted only once:
%CONT>>
